\documentclass[12pt]{article}

\usepackage{amsmath}
\usepackage{amssymb}
\usepackage{physics}
\usepackage{hyperref}

\title{Parallelism in Affine Space}
\author{}
\date{}

\begin{document}

\maketitle

\section{Parallelism in Affine Space}

\begin{definition}[Parallel affine subspaces]
Let
\[
S(x,U)
\qquad\text{and}\qquad
S(x',W)
\]
be affine subspaces of \(\mathcal A\). They are said to be parallel if and only if
they have the same direction space; that is,
\[
S(x,U)\parallel S(x',W)
\quad\Longleftrightarrow\quad
U=W.
\]
\end{definition}

\begin{proposition}
Let \(S(x,U)\) be an affine subspace of \(\mathcal A\). For every
\(\mathbf v\in V\),
\[
T_{\mathbf v}\bigl(S(x,U)\bigr)
=
S\bigl(\mathbf v(x),U\bigr).
\]
Consequently, every translate of an affine subspace is parallel to the
original affine subspace.
\end{proposition}

\begin{proof}
By definition,
\[
T_{\mathbf v}\bigl(S(x,U)\bigr)
=
\left\{
y\in \mathcal A
\;\middle|\;
y=\mathbf v\bigl(\mathbf u(x)\bigr)
\text{ for some }\mathbf u\in U
\right\}.
\]

Let
\[
y\in T_{\mathbf v}\bigl(S(x,U)\bigr).
\]
Then there exists \(\mathbf u\in U\) such that
\[
y=\mathbf v\bigl(\mathbf u(x)\bigr).
\]
Using the properties of the action of \(V\) on \(\mathcal A\), we obtain
\[
y
=
(\mathbf v+\mathbf u)(x)
=
\mathbf u\bigl(\mathbf v(x)\bigr).
\]
Therefore
\[
y\in S\bigl(\mathbf v(x),U\bigr),
\]
and hence
\[
T_{\mathbf v}\bigl(S(x,U)\bigr)
\subseteq
S\bigl(\mathbf v(x),U\bigr).
\]

Conversely, let
\[
y\in S\bigl(\mathbf v(x),U\bigr).
\]
Then there exists \(\mathbf u\in U\) such that
\[
y=\mathbf u\bigl(\mathbf v(x)\bigr).
\]
Again using the properties of the action, we obtain
\[
y
=
(\mathbf u+\mathbf v)(x)
=
\mathbf v\bigl(\mathbf u(x)\bigr).
\]
Since
\[
\mathbf u(x)\in S(x,U),
\]
it follows that
\[
y\in T_{\mathbf v}\bigl(S(x,U)\bigr).
\]
Therefore
\[
S\bigl(\mathbf v(x),U\bigr)
\subseteq
T_{\mathbf v}\bigl(S(x,U)\bigr).
\]

We conclude that
\[
T_{\mathbf v}\bigl(S(x,U)\bigr)
=
S\bigl(\mathbf v(x),U\bigr).
\]

Since both affine subspaces have direction space \(U\),
\[
T_{\mathbf v}\bigl(S(x,U)\bigr)
\parallel
S(x,U).
\]
\end{proof}

\begin{proposition}
Let
\[
S=S(x,U)
\qquad\text{and}\qquad
S'=S(y,W)
\]
be affine subspaces of \(\mathcal A\).

If \(S\parallel S'\), then either
\[
S=S'
\]
or
\[
S\cap S'=\varnothing.
\]

Moreover, if \(S\) and \(S'\) are hyperplanes, then
\[
S\parallel S'
\quad\Longleftrightarrow\quad
S=S'
\ \text{or}\
S\cap S'=\varnothing.
\]
\end{proposition}

\begin{proof}
Suppose first that \(S\parallel S'\). By definition, their direction
spaces coincide, so
\[
U=W.
\]

Assume that
\[
S\cap S'\neq\varnothing,
\]
and choose a point \(z\in S\cap S'\). Since \(z\in S\) and
\(z\in S'\), there exist vectors \(\mathbf b,\mathbf c\in U\) such
that
\[
z=\mathbf b(x)=\mathbf c(y).
\]
It follows that
\[
y=(-\mathbf c)(z)
  =(-\mathbf c)\bigl(\mathbf b(x)\bigr)
  =(\mathbf b-\mathbf c)(x).
\]
Since \(\mathbf b-\mathbf c\in U\), we have \(y\in S\).

Now let \(p\in S'\). There exists \(\mathbf u\in U\) such that
\[
p=\mathbf u(y).
\]
Therefore,
\[
p
 =\mathbf u\bigl((\mathbf b-\mathbf c)(x)\bigr)
 =(\mathbf u+\mathbf b-\mathbf c)(x).
\]
Because
\[
\mathbf u+\mathbf b-\mathbf c\in U,
\]
we obtain \(p\in S\). Hence
\[
S'\subseteq S.
\]
Interchanging \(S\) and \(S'\) gives
\[
S\subseteq S',
\]
and consequently
\[
S=S'.
\]
Thus, if \(S\) and \(S'\) are parallel, then either
\[
S=S'
\]
or
\[
S\cap S'=\varnothing.
\]

We now assume that \(S\) and \(S'\) are hyperplanes and prove the
converse. If \(S=S'\), then their direction spaces are equal, and
therefore \(S\parallel S'\).

It remains to prove that
\[
S\cap S'=\varnothing
\quad\Longrightarrow\quad
S\parallel S'.
\]
Equivalently, we prove the contrapositive
\[
S\not\parallel S'
\quad\Longrightarrow\quad
S\cap S'\neq\varnothing.
\]

Suppose, therefore, that \(S\not\parallel S'\). Then
\[
U\neq W.
\]
Since \(U\) and \(W\) are distinct \((n-1)\)-dimensional linear
subspaces of \(V\), we have
\[
U+W=V.
\]

Let
\[
\mathbf d=\overrightarrow{xy}.
\]
Since \(\mathbf d\in V=U+W\), there exist
\[
\mathbf b\in U
\qquad\text{and}\qquad
\mathbf c\in W
\]
such that
\[
\mathbf d=\mathbf b+\mathbf c.
\]
Consequently,
\[
y=\mathbf d(x)=(\mathbf b+\mathbf c)(x),
\]
and hence
\[
(-\mathbf c)(y)=\mathbf b(x). 
\]

The point \(\mathbf b(x)\) belongs to \(S\), since \(\mathbf b\in U\),
while the same point, written as \((-\mathbf c)(y)\), belongs to
\(S'\), since \(-\mathbf c\in W\). Therefore,
\[
S\cap S'\neq\varnothing.
\]

This proves the contrapositive and hence
\[
S\cap S'=\varnothing
\quad\Longrightarrow\quad
S\parallel S'.
\]
The proof is complete.
\end{proof}

\begin{proposition}[The fifth parallel axiom]
Let \(S\) be a \(r\)-dimensional affine subspace of \(\mathcal A\), and let
\(x\in \mathcal A\). Then there exists one and only one \(r\)-dimensional
affine subspace \(S'\) of \(\mathcal A\) which contains \(x\) and is parallel
to \(S\).
\end{proposition}

\begin{proof}
Let
\[
S=S(y,U),
\]
where \(U\) is a \(r\)-dimensional linear subspace of \(V\). Define
\[
\mathbf d=\overrightarrow{yx}.
\]
Then
\[
T_{\mathbf d}(y)=x.
\]
We therefore define
\[
S'
   =S\bigl(T_{\mathbf d}(y),U\bigr)
   =S(x,U).
\]
It follows immediately that
\[
x\in S'
\qquad\text{and}\qquad
S'\parallel S,
\]
since \(S\) and \(S'\) have the same direction space \(U\). Moreover,
\[
\dim S'=\dim U=r.
\]
This proves existence.

To prove uniqueness, let \(S''\) be a \(r\)-dimensional affine
subspace containing \(x\) and parallel to \(S\). Since \(S''\parallel
S\), its direction space is \(U\). Since \(x\in S''\), we may take
\(x\) as its reference point, and therefore
\[
S''=S(x,U)=S'.
\]
Hence \(S'\) is unique.
\end{proof}

\end{document}
